\documentclass[12pt, a4paper]{article}

\usepackage[utf8]{inputenc}
\usepackage[T1]{fontenc}
\usepackage[brazil]{babel}
\usepackage{amsmath}
\usepackage[margin=2.5cm]{geometry}
\usepackage{hyperref}
\usepackage{enumitem}
\usepackage{titlesec}
\usepackage{booktabs}
\usepackage{array}
\usepackage{xcolor}
\usepackage{parskip}

\definecolor{primary}{HTML}{1E3A5F}
\definecolor{accent}{HTML}{2E86AB}
\definecolor{lightgray}{HTML}{F5F5F5}

\titleformat{\section}{\large\bfseries\color{primary}}{}{0em}{}[\titlerule]
\titleformat{\subsection}{\normalsize\bfseries\color{accent}}{}{0em}{}

\hypersetup{
    colorlinks=true,
    linkcolor=accent,
    urlcolor=accent
}

\title{
    \vspace{-1cm}
    {\large\color{accent} Documento de Referência Técnica} \\
    \vspace{0.3cm}
    {\LARGE\bfseries\color{primary} Alvenar — MVP} \\
    \vspace{0.3cm}
    {\normalsize Decisões técnicas e escopo acordados até o momento}
}
\author{Equipe Alvenar}
\date{\today}

\begin{document}

\maketitle
\vspace{0.5cm}
\hrule
\vspace{1cm}

% -----------------------------------------
\section{Stack Tecnológica}

\subsection{Frontend}

\begin{itemize}[leftmargin=1.5em]
    \item \textbf{React + TypeScript} — interface web do síndico e painel de transparência do morador.
    A tipagem estática do TypeScript reduz erros em formulários complexos (cadastro de equipamentos,
    upload de laudos) e facilita a manutenção à medida que o produto cresce.
    \item \textbf{Tailwind CSS} — estilização utilitária inline, sem arquivos CSS separados.
    Elimina a necessidade de nomear classes e manter consistência manual de espaçamento,
    especialmente vantajoso para um time enxuto.
    \item \textbf{shadcn/ui} — biblioteca de componentes cujo código é copiado diretamente
    para o projeto, sem dependência de versão externa. Fornece prontos os componentes mais
    usados no Alvenar: \texttt{Table}, \texttt{Badge}, \texttt{Dialog}, \texttt{Card},
    \texttt{Form} e \texttt{Progress}.
    \item \textbf{Recharts} — gráficos (barras empilhadas do calendário de inspeções e
    barra de score de saúde), com integração nativa ao shadcn/ui.
    \item \textbf{react-hook-form + Zod} — gerenciamento e validação de formulários.
    Os schemas Zod podem ser compartilhados entre frontend e backend, garantindo
    consistência de tipos de ponta a ponta.
\end{itemize}

\subsection{Backend}

\begin{itemize}[leftmargin=1.5em]
    \item \textbf{Node.js + Fastify} — servidor HTTP responsável pelas regras de negócio,
    agendamento de alertas e controle de acesso. O Fastify foi escolhido sobre o Express por
    oferecer suporte nativo a TypeScript, melhor performance e validação de schema integrada.
    \item \textbf{Prisma ORM} — camada de acesso ao banco de dados. Oferece migrations
    versionadas, type-safety de ponta a ponta com TypeScript e simplifica consultas relacionais.
    \item \textbf{node-cron} — agendamento do job diário de alertas (executa às 8h,
    horário de Brasília).
    \item \textbf{Nodemailer + Gmail OAuth2} — envio de e-mails transacionais. A autenticação
    via OAuth2 elimina o uso de senhas de app e é mais robusta para produção.
\end{itemize}

\subsection{Banco de Dados}

\begin{itemize}[leftmargin=1.5em]
    \item \textbf{PostgreSQL} — banco relacional escolhido por afinidade com a estrutura dos
    dados do domínio (detalhado na Seção~\ref{sec:banco}).
\end{itemize}

\subsection{Futuro — Fora do MVP}

\begin{itemize}[leftmargin=1.5em]
    \item \textbf{Flutter} — aplicativo mobile para síndicos e moradores. Presente na stack
    final, mas não no MVP.
\end{itemize}

\vspace{0.5cm}

% -----------------------------------------
\section{Justificativa: PostgreSQL vs.\ NoSQL}
\label{sec:banco}

Os dados do Alvenar são \textbf{altamente relacionais} por natureza:

\[
\text{Condomínio} \rightarrow \text{Equipamentos} \rightarrow \text{Vistorias} \rightarrow \text{Documentos/Laudos}
\]

As razões para manter o PostgreSQL no MVP são:

\begin{enumerate}[leftmargin=1.5em]
    \item \textbf{Integridade referencial} — chaves estrangeiras garantem que nenhuma vistoria
    fique órfã de equipamento, e nenhum laudo seja vinculado a uma vistoria inexistente.
    \item \textbf{Consistência transacional (ACID)} — crítico ao registrar uma vistoria junto
    com seu documento em uma única operação atômica.
    \item \textbf{Queries relacionais} — consultas como ``todos os equipamentos vencidos de um
    condomínio'' são triviais em SQL e custosas em documentos NoSQL.
    \item \textbf{Prisma + Postgres} — integração madura, com suporte completo a migrations e
    geração automática de tipos TypeScript.
\end{enumerate}

\textbf{Quando NoSQL faria sentido?} Apenas se houvesse necessidade de armazenar dados
não estruturados em larga escala (ex.: logs de IoT de sensores de elevador) — o que está
fora do escopo do MVP.

\vspace{0.5cm}

% -----------------------------------------
\section{Casos de Uso Essenciais do MVP}

O critério de inclusão é direto: apenas o que valida a hipótese central do produto.

\begin{quote}
\textit{``A automação de alertas e a centralização digital de documentos reduz
o esquecimento de manutenções críticas.''}
\end{quote}

\subsection{UC-01 — Gestão de Vistorias}

\begin{itemize}[leftmargin=1.5em]
    \item Cadastrar equipamentos obrigatórios (elevador, para-raios, caixa d'água etc.)
    \item Definir periodicidade de cada manutenção conforme NBR 5674
    \item Visualizar cronograma de vencimentos
\end{itemize}

\subsection{UC-02 — Alertas Proativos}

\begin{itemize}[leftmargin=1.5em]
    \item Receber notificação antes do vencimento de um prazo técnico
    \item Marcar vistoria como realizada
    \item Registrar responsável e observações da execução
\end{itemize}

\subsection{UC-03 — Repositório de Documentos}

\begin{itemize}[leftmargin=1.5em]
    \item Anexar laudos e certificados a cada vistoria
    \item Acessar histórico técnico de um equipamento
    \item Base da \textit{memória técnica auditável} — principal diferencial do produto
\end{itemize}

\subsection{UC-04 — Painel de Saúde do Edifício}

\begin{itemize}[leftmargin=1.5em]
    \item Síndico visualiza status geral: \textcolor{green!60!black}{\textbf{em dia}} /
    \textcolor{orange}{\textbf{vencendo}} / \textcolor{red}{\textbf{vencido}}
    \item Morador acessa visão simplificada de transparência (somente leitura)
\end{itemize}

\vspace{0.5cm}

% -----------------------------------------
\section{Fora do Escopo do MVP}

Os itens abaixo são relevantes para versões futuras, mas \textbf{não testam a hipótese
central} e foram deliberadamente excluídos:

\begin{itemize}[leftmargin=1.5em]
    \item Score de saúde com Inteligência Artificial
    \item Validação de autenticidade de laudos (ex.: integração com cartório digital)
    \item Integração com ERPs condominiais existentes (ex.: Superlógica, Condomínio21)
    \item Aplicativo mobile (Flutter) — presente na stack final, mas não no MVP
    \item \textbf{Sugestão de corretiva pelo morador} — moradores poderiam abrir chamados sugerindo itens para manutenção corretiva, a serem aprovados ou ignorados pelo síndico
    \item \textbf{Controle de acesso a documentos} — arquivos em \texttt{/uploads} devem ser protegidos por middleware de autenticação, garantindo que apenas usuários vinculados ao condomínio consigam acessar os laudos. Hoje o servidor de arquivos estáticos serve arquivos publicamente
\end{itemize}

\vspace{0.5cm}

% -----------------------------------------
\section{Arquitetura de Software}
\label{sec:arquitetura}

Adotada uma \textbf{arquitetura em camadas simples}, priorizada pela manutenibilidade
por um time enxuto (um engenheiro de software + analistas/suporte):

\begin{verbatim}
src/
+-- routes/          # Roteamento Fastify
+-- controllers/     # Recebe req, chama service, devolve res
+-- services/        # Regras de negócio
+-- repositories/    # Acesso ao banco via Prisma
+-- jobs/            # Cron de alertas (node-cron)
+-- middlewares/     # Autenticação, error handler global
\end{verbatim}

\textbf{Princípios adotados:}

\begin{itemize}[leftmargin=1.5em]
    \item Controllers são \textit{burros de propósito} — nenhuma regra de negócio vive ali.
    \item Services contêm toda a lógica do domínio e são testáveis de forma isolada.
    \item Repositories encapsulam o Prisma — facilitam substituição futura sem alterar services.
    \item Um error handler global no Fastify elimina try/catch repetitivo nos controllers.
\end{itemize}

\textbf{Decisões deliberadamente evitadas no MVP:}

\begin{itemize}[leftmargin=1.5em]
    \item Microserviços — complexidade operacional desnecessária nesta fase.
    \item Clean Architecture / Hexagonal completo — overhead burocrático incompatível com
    time pequeno e prazo de MVP.
    \item Feature folders — organização por camada é suficiente até o time crescer.
\end{itemize}

\vspace{0.5cm}

% -----------------------------------------
\section{Sistema de Alertas}
\label{sec:alertas}

\subsection{Fluxo de Disparo}

O job de alertas roda diariamente às \textbf{8h (America/Sao\_Paulo)} via \texttt{node-cron}
e processa todas as inspeções com status \texttt{SCHEDULED} ou \texttt{OVERDUE}.

Para cada inspeção, há três gatilhos independentes:

\begin{center}
\begin{tabular}{lll}
\toprule
\textbf{Gatilho} & \textbf{Condição} & \textbf{Frequência} \\
\midrule
\texttt{DAYS\_15\_BEFORE} & 15 dias antes do \texttt{scheduledDate} & Uma vez \\
\texttt{DAYS\_7\_BEFORE}  & 7 dias antes do \texttt{scheduledDate}  & Uma vez \\
\texttt{DAILY\_REMINDER}  & Últimos 7 dias E status $\neq$ \texttt{DONE} & Diário \\
\bottomrule
\end{tabular}
\end{center}

\subsection{Destinatários}

Todos os usuários com papel \texttt{MANAGER} vinculados ao condomínio do equipamento
recebem o alerta.

\subsection{Idempotência}

Antes de enviar, o sistema verifica se já existe um \texttt{Alert} com o mesmo
\texttt{inspectionId} + \texttt{trigger} + \texttt{status: SENT} no mesmo dia.
Isso previne duplicidade em caso de reexecução do job.

\subsection{Tratamento de Falhas}

Falhas de envio são registradas no model \texttt{Alert} com \texttt{status: FAILED}
e a mensagem de erro. O job continua processando os demais alertas — sem retry
automático no MVP.

\subsection{Edge Cases Documentados}

\begin{itemize}[leftmargin=1.5em]
    \item \textbf{Servidor fora no dia do gatilho} — o alerta D-15 ou D-7 não será
    reenviado no dia seguinte. Mitigação simples: ampliar a janela de verificação
    para $\pm$1 dia (\texttt{daysUntilDue <= 15 \&\& >= 14}).
    \item \textbf{Inspeção marcada como DONE durante a semana final} — o cron para
    de enviar automaticamente, pois o \texttt{where} filtra apenas \texttt{SCHEDULED}
    e \texttt{OVERDUE}. Depende da atualização correta de status ao registrar execução.
\end{itemize}

\vspace{0.5cm}

% -----------------------------------------
\section{Modelagem do Banco de Dados}
\label{sec:modelagem}

\subsection{Decisões de Modelagem}

As seguintes premissas guiaram o design do schema:

\begin{itemize}[leftmargin=1.5em]
    \item \textbf{Síndico profissional (N:N)} — um usuário pode gerenciar múltiplos condomínios
    e um condomínio pode ter múltiplos usuários. Resolvido via tabela de junção
    \texttt{UserCondominium} com campo \texttt{role}.
    \item \textbf{Dois papéis apenas} — \texttt{MANAGER} (síndico/zelador, front de gestão)
    e \texttt{RESIDENT} (somente leitura, front de transparência). Papel de zelador
    diferenciado entra em versão futura.
    \item \textbf{Planos de manutenção baseados na NBR 5674} — pré-definidos pelo sistema
    mas editáveis pelo síndico. Flag \texttt{isCustomized} registra se o padrão foi alterado.
    \item \textbf{Três tipos de manutenção} — rotina, preventiva e corretiva. A corretiva
    é registrada manualmente pelo síndico após uma falha.
    \item \textbf{Histórico de alertas persistido} — cada alerta enviado é registrado com
    \texttt{trigger}, \texttt{status} e \texttt{timestamp} para rastreabilidade, auditoria
    e idempotência.
    \item \textbf{Documentos desacoplados} — laudos vinculados à inspeção, não ao equipamento,
    pois cada ciclo gera laudos distintos.
    \item \textbf{Armazenamento local no MVP} — arquivos servidos com path estruturado:
    \texttt{/uploads/inspections/\{inspectionId\}/\{fileName\}}.
    Facilita migração futura para S3 sem alterar o schema.
    \item \textbf{Canal de alerta} — apenas \texttt{EMAIL} no MVP via Nodemailer + Gmail OAuth2.
    Push notification entra com o app mobile Flutter em versão futura.
\end{itemize}

\subsection{Diagrama de Entidades (ER simplificado)}

\begin{verbatim}
  User ---------------- UserCondominium ---- Condominium
    |                       (role)                |
    |                                             |
    +-- registeredBy ---- Inspection ---- MaintenancePlan -- Equipment
                          /        \
                     Document     Alert
                                 (trigger)
\end{verbatim}

\subsection{Schema Prisma}

\begin{verbatim}
// -- Enums --------------------------------------------------

enum Role {
  MANAGER   // síndico — acesso total (front de gestão)
  RESIDENT  // morador — somente leitura (front de transparência)
}

enum EquipmentType {
  ELEVATOR
  LIGHTNING_ROD
  WATER_TANK
  OTHER
}

enum MaintenanceType {
  ROUTINE     // rotina — periodicidade fixa
  PREVENTIVE  // preventiva — programada
  CORRECTIVE  // corretiva — após falha, registrada manualmente
}

enum InspectionStatus {
  SCHEDULED  // agendada
  DONE       // realizada
  OVERDUE    // vencida
}

enum AlertStatus {
  PENDING
  SENT
  FAILED
}

enum AlertChannel {
  EMAIL
}

enum AlertTrigger {
  DAYS_15_BEFORE  // disparo único em D-15
  DAYS_7_BEFORE   // disparo único em D-7
  DAILY_REMINDER  // reenvio diário na semana final
}

// -- Models -------------------------------------------------

model User {
  id           String            @id @default(uuid())
  name         String
  email        String            @unique
  passwordHash String
  createdAt    DateTime          @default(now())
  updatedAt    DateTime          @updatedAt

  condominiums UserCondominium[]
  inspections  Inspection[]
}

model Condominium {
  id           String            @id @default(uuid())
  name         String
  cnpj         String?           @unique
  address      String
  createdAt    DateTime          @default(now())
  updatedAt    DateTime          @updatedAt

  users        UserCondominium[]
  equipment    Equipment[]
}

model UserCondominium {
  id            String      @id @default(uuid())
  userId        String
  condominiumId String
  role          Role
  createdAt     DateTime    @default(now())

  user          User        @relation(fields: [userId],        references: [id])
  condominium   Condominium @relation(fields: [condominiumId], references: [id])

  @@unique([userId, condominiumId])
}

model Equipment {
  id            String            @id @default(uuid())
  condominiumId String
  name          String
  type          EquipmentType
  createdAt     DateTime          @default(now())
  updatedAt     DateTime          @updatedAt

  condominium      Condominium      @relation(fields: [condominiumId], references: [id])
  maintenancePlans MaintenancePlan[]
}

model MaintenancePlan {
  id               String          @id @default(uuid())
  equipmentId      String
  type             MaintenanceType
  frequencyDays    Int
  alertDaysBeforeDue Int[]         @default([15, 7])
  description      String
  isCustomized     Boolean         @default(false)
  createdAt        DateTime        @default(now())
  updatedAt        DateTime        @updatedAt

  equipment     Equipment       @relation(fields: [equipmentId], references: [id])
  inspections   Inspection[]
}

model Inspection {
  id                String           @id @default(uuid())
  maintenancePlanId String
  scheduledById     String
  executedById      String?
  scheduledDate     DateTime
  executedAt        DateTime?
  status            InspectionStatus
  notes             String?
  createdAt         DateTime         @default(now())
  updatedAt         DateTime         @updatedAt

  maintenancePlan   MaintenancePlan  @relation(fields: [maintenancePlanId], references: [id])
  scheduledBy       User             @relation("ScheduledBy", fields: [scheduledById], references: [id])
  executedBy        User?            @relation("ExecutedBy",  fields: [executedById],  references: [id])
  documents         Document[]
  alerts            Alert[]
}

model Document {
  id           String     @id @default(uuid())
  inspectionId String
  fileName     String
  url          String
  mimeType     String
  sizeBytes    Int
  createdAt    DateTime   @default(now())

  inspection   Inspection @relation(fields: [inspectionId], references: [id])
}

model Alert {
  id           String        @id @default(uuid())
  inspectionId String
  channel      AlertChannel
  trigger      AlertTrigger
  status       AlertStatus   @default(PENDING)
  sentAt       DateTime?
  error        String?
  createdAt    DateTime      @default(now())

  inspection   Inspection    @relation(fields: [inspectionId], references: [id])
}
\end{verbatim}

\subsection{Notas de Implementação}

\begin{itemize}[leftmargin=1.5em]
    \item \textbf{Status calculado vs.\ persistido} — \texttt{InspectionStatus} pode ser
    calculado em runtime comparando \texttt{scheduledDate} com \texttt{NOW()}, mas
    persistir facilita queries do painel. Recomenda-se um job cron que atualiza o campo
    periodicamente (o mesmo job de alertas pode acumular essa responsabilidade).
    \item \textbf{scheduledById vs executedById} — separação explícita entre quem agendou
    e quem executou a inspeção. \texttt{executedById} é nullable até a inspeção ser concluída.
    \item \textbf{Corretiva} — \texttt{scheduledDate} recebe a data do registro da falha.
    \texttt{frequencyDays} no \texttt{MaintenancePlan} é ignorado para este tipo.
    \item \textbf{Índices recomendados} — \texttt{Inspection.scheduledDate} (para o cron
    de alertas) e \texttt{Equipment.condominiumId} (para o painel de saúde).
    \item \textbf{Soft delete} — adicionar \texttt{deletedAt DateTime?} nas entidades
    principais é \textbf{obrigatório} dado que a memória técnica auditável é o principal
    diferencial do produto. Remoção física de registros contradiz esse valor.
\end{itemize}

\vspace{0.5cm}

% -----------------------------------------
\section{Componentes do Frontend}
\label{sec:frontend}

Os componentes estão organizados por página. Cada componente lista a query Prisma
correspondente que alimenta seus dados.

\subsection{Autenticação — \texttt{/login} · \texttt{/register}}

\begin{itemize}[leftmargin=1.5em]
    \item \textbf{LoginForm} — \texttt{User.findUnique(\{ email \})}
    \item \textbf{RegisterForm} — \texttt{User.create()}
\end{itemize}

\subsection{Dashboard — \texttt{/dashboard} (MANAGER)}

\begin{itemize}[leftmargin=1.5em]
    \item \textbf{HealthScoreCard} — score e status geral do condomínio \\
          Query: \texttt{Inspection.count(\{ condominiumId, status \})}
    \item \textbf{MetricsSummary} — contagem de inspeções em dia / vencendo / vencidas \\
          Query: \texttt{Inspection.groupBy(status)}
    \item \textbf{EquipmentList} — lista resumida com badge de status por equipamento \\
          Query: \texttt{Equipment.findMany(\{ condominiumId, include: nextInspection \})}
    \item \textbf{AlertFeed} — alertas recentes do condomínio \\
          Query: \texttt{Alert.findMany(\{ condominiumId, orderBy: createdAt, take: 10 \})}
    \item \textbf{InspectionCalendar} — inspeções agendadas nos próximos 30 dias \\
          Query: \texttt{Inspection.findMany(\{ scheduledDate: \{ gte: today, lte: +30d \} \})}
\end{itemize}

\textbf{Nota:} O \texttt{EquipmentList} do dashboard busca apenas o próximo vencimento
por equipamento — endpoint separado do \texttt{EquipmentTable} da página de equipamentos,
que traz todos os planos de manutenção.

\subsection{Equipamentos — \texttt{/equipment}}

\begin{itemize}[leftmargin=1.5em]
    \item \textbf{EquipmentTable} — tabela completa de equipamentos com planos \\
          Query: \texttt{Equipment.findMany(\{ condominiumId, include: maintenancePlans \})}
    \item \textbf{AddEquipmentForm} — cadastro de equipamento + plano de manutenção NBR \\
          Queries: \texttt{Equipment.create()} + \texttt{MaintenancePlan.create()}
\end{itemize}

\subsection{Detalhe do Equipamento — \texttt{/equipment/:id}}

\begin{itemize}[leftmargin=1.5em]
    \item \textbf{EquipmentHeader} — nome, tipo e status atual \\
          Query: \texttt{Equipment.findUnique(\{ id, include: maintenancePlans \})}
    \item \textbf{InspectionHistory} — histórico de vistorias com documentos vinculados \\
          Query: \texttt{Inspection.findMany(\{ equipmentId, include: documents, orderBy: scheduledDate \})}
    \item \textbf{DocumentList} — laudos e certificados de uma inspeção \\
          Query: \texttt{Document.findMany(\{ inspectionId \})}
    \item \textbf{MarkDoneForm} — registrar execução e fazer upload de laudo \\
          Queries: \texttt{Inspection.update(\{ status: DONE, executedById, executedAt \})} + \texttt{Document.create(\{ url, mimeType \})}
\end{itemize}

\subsection{Painel do Morador — \texttt{/resident} (RESIDENT)}

\begin{itemize}[leftmargin=1.5em]
    \item \textbf{BuildingHealth} — score de saúde em modo somente leitura \\
          Query: \texttt{Inspection.count(\{ condominiumId, status \})}
    \item \textbf{PublicEquipmentList} — status dos equipamentos sem dados sensíveis \\
          Query: \texttt{Equipment.findMany(\{ condominiumId, select: \{ name, type, status \} \})}
    \item \textbf{LastInspectionBadge} — data da última vistoria realizada por equipamento \\
          Query: \texttt{Inspection.findFirst(\{ equipmentId, status: DONE, orderBy: executedAt desc \})}
\end{itemize}

\textbf{Nota de segurança:} O painel do morador usa \texttt{select} explícito no Prisma
para não expor campos sensíveis como \texttt{scheduledById}, \texttt{executedById} ou
\texttt{notes}. Esse filtro deve ser garantido na camada de serviço, não apenas no frontend.

\vspace{0.5cm}

% -----------------------------------------
\section{Dashboard — Layout e Decisões de UX}
\label{sec:dashboard}

\subsection{Wireframe — Visão do Síndico (MANAGER)}

\begin{verbatim}
+-------------------------------------------------------------+
|  Alvenar            [Cond. Parque das Flores]    MANAGER v  |
+-------------------------------------------------------------+
|                                                             |
|  +-----------+  +-----------+  +-----------+  +---------+ |
|  |Equipament.|  |  Em dia   |  | Vencendo  |  |Vencidos | |
|  |    12     |  |     7     |  |     3     |  |    2    | |
|  +-----------+  +-----------+  +-----------+  +---------+ |
|                                                             |
|  +-----------------------------------------------------+   |
|  |  Saúde do condomínio                       72 pts   |   |
|  |  Em dia    ############...................   7      |   |
|  |  Vencendo  #####..........................   3      |   |
|  |  Vencidos  ###............................   2      |   |
|  |  [====================|..........] Regular  |   |
|  +-----------------------------------------------------+   |
|                                                             |
|  +--------------------------+ +--------------------------+ |
|  | Equipamentos             | | Alertas recentes          | |
|  | Elevador A   [vencendo]  | | * Para-raios — vencido    | |
|  | Para-raios   [vencido]   | | * Bomba — vencida 38d     | |
|  | Caixa d'água [em dia]    | | * Elevador A — D-15 env.  | |
|  | Elevador B   [em dia]    | | * Caixa d'água — concluído| |
|  | Bomba recalq.[vencido]   | |                           | |
|  | [Ver todos ->]            | | [Ver histórico ->]         | |
|  +--------------------------+ +--------------------------+ |
|                                                             |
|  +-----------------------------------------------------+   |
|  |  Inspeções nos próximos 30 dias                     |   |
|  |  # Realizadas  # Agendadas  # Vencidas              |   |
|  |  |   |   |   |   |   |   |   |   |   |             |   |
|  |  15  18  21  24  27  30  02  05  08  11             |   |
|  +-----------------------------------------------------+   |
+-------------------------------------------------------------+
\end{verbatim}

\subsection{Wireframe — Visão do Morador (RESIDENT)}

\begin{verbatim}
+-------------------------------------------------------------+
|  Alvenar            [Cond. Parque das Flores]    RESIDENT v |
+-------------------------------------------------------------+
|                                                             |
|  +-----------------------------------------------------+   |
|  |  Saúde do edifício                                  |   |
|  |                                                     |   |
|  |         Status geral: Regular (72 pts)              |   |
|  |  [==================|................]              |   |
|  +-----------------------------------------------------+   |
|                                                             |
|  +-----------------------------------------------------+   |
|  |  Equipamentos do condomínio                         |   |
|  |                                                     |   |
|  |  Elevador Social A        [vencendo]                |   |
|  |  Para-raios               [vencido]                 |   |
|  |  Caixa d'água             [em dia]    últ: mar/26   |   |
|  |  Elevador Social B        [em dia]    últ: fev/26   |   |
|  |  Bomba de recalque        [vencido]                 |   |
|  |                                                     |   |
|  |  * Datas de execução exibidas apenas quando         |   |
|  |    a vistoria estiver concluída (status DONE)       |   |
|  +-----------------------------------------------------+   |
+-------------------------------------------------------------+
\end{verbatim}

\subsection{Decisões de UX}

\textbf{Score de saúde numérico vs.\ qualitativo.}
O dashboard exibe o score numérico (0--100) para o síndico, acompanhado de um
rótulo qualitativo (Crítico / Regular / Bom / Excelente). O morador vê apenas
o rótulo qualitativo, reduzindo a sobrecarga de informação para quem não precisa
do detalhe técnico. A fórmula base é:

\[
\text{Score} = \frac{\text{inspeções em dia}}{\text{total de inspeções}} \times 100
\]

Versões futuras podem ponderar por criticidade do equipamento (elevador
tem peso maior que limpeza de calha).

\textbf{Separação de visões por papel.}
O painel do morador é uma tela distinta, não um modo de leitura do painel do
síndico. Isso evita lógica condicional complexa no frontend e garante que campos
sensíveis (responsável pela execução, observações internas, datas de agendamento)
nunca apareçam acidentalmente para residentes.

\textbf{Dados exibidos ao morador.}
Apenas nome do equipamento, tipo e status são exibidos publicamente. A data da
última vistoria concluída (\texttt{executedAt}) é exibida quando disponível —
reforça a transparência sem expor o planejamento interno do síndico.

\textbf{Ações rápidas no dashboard.}
O botão ``Marcar como realizada'' aparece diretamente na lista de equipamentos
vencidos ou vencendo, sem necessidade de navegar até a página de detalhe.
Reduz o atrito no fluxo mais frequente do síndico.

\textbf{Gráfico de inspeções.}
O gráfico de barras empilhadas cobre os próximos 30 dias e usa as três cores
semânticas do sistema (verde / âmbar / vermelho). Alimentado diretamente pela
query \texttt{Inspection.findMany} filtrada por \texttt{scheduledDate}.


\vspace{0.5cm}

% -----------------------------------------
\section{Arquitetura do Frontend}
\label{sec:frontarq}

\subsection{Stack Definida}

\begin{itemize}[leftmargin=1.5em]
    \item \textbf{React Router v6} --- roteamento. Preferido ao TanStack Router
    pela maturidade e ecossistema maior, sem custo adicional para o MVP.
    \item \textbf{TanStack Query} --- gerenciamento de server state (fetch, cache,
    loading, erro e refetch). Suficiente para o Alvenar, onde todo estado relevante
    vem do servidor.
    \item \textbf{Context API} --- estado global client-side limitado ao usuário
    autenticado e ao condomínio ativo (\texttt{AuthContext}).
    \item \textbf{react-hook-form + Zod} --- formulários e validação. Os schemas Zod
    são compartilhados com o backend, garantindo consistência de tipos de ponta a ponta.
    \item \textbf{Recharts} --- gráficos, com integração nativa ao shadcn/ui.
\end{itemize}

\subsection{Estrutura de Pastas}

\begin{verbatim}
src/
+-- components/
|   +-- ui/                  # shadcn (Button, Card, Badge, Dialog...)
|   +-- HealthScoreCard.tsx
|   +-- EquipmentTable.tsx
|   +-- AlertFeed.tsx
|   +-- InspectionCalendar.tsx
|
+-- pages/
|   +-- dashboard/
|   |   +-- DashboardPage.tsx
|   +-- equipment/
|   |   +-- EquipmentListPage.tsx
|   |   +-- EquipmentDetailPage.tsx
|   +-- resident/
|   |   +-- ResidentPage.tsx
|   +-- auth/
|       +-- LoginPage.tsx
|       +-- RegisterPage.tsx
|
+-- hooks/
|   +-- useEquipment.ts      # TanStack Query wrappers
|   +-- useInspections.ts
|   +-- useAlerts.ts
|
+-- services/
|   +-- api.ts               # instancia axios com base URL + auth header
|   +-- equipment.service.ts
|   +-- inspection.service.ts
|
+-- context/
|   +-- AuthContext.tsx      # user, condominiumId, role
|
+-- types/
|   +-- index.ts             # Equipment, Inspection, Alert, User...
|
+-- lib/
|   +-- schemas.ts           # Zod schemas compartilhados com o backend
|   +-- utils.ts             # formatadores de data, status helpers
|
+-- App.tsx                  # rotas React Router
\end{verbatim}

\subsection{AuthContext}

O \texttt{AuthContext} é o único estado global client-side do MVP. Armazena
o usuário autenticado, o condomínio ativo e o papel (\texttt{role}), tornando
esses dados acessíveis em qualquer componente sem prop drilling.

\begin{verbatim}
// context/AuthContext.tsx

interface AuthState {
  user: { id: string; name: string; email: string } | null
  condominiumId: string | null
  role: 'MANAGER' | 'RESIDENT' | null
  token: string | null
}

const AuthContext = createContext<AuthState & {
  login: (token: string, user: AuthState['user'],
          condominiumId: string, role: AuthState['role']) => void
  logout: () => void
}>(null!)
\end{verbatim}

O token JWT é armazenado em memória (não em \texttt{localStorage}) e
re-hidratado via refresh token em cookie \texttt{httpOnly} ao recarregar
a página --- padrão mais seguro contra ataques XSS.

\subsection{Hooks com TanStack Query}

Cada entidade do domínio tem um hook dedicado que encapsula a query e
expõe os dados já tipados. Exemplo para equipamentos:

\begin{verbatim}
// hooks/useEquipment.ts

export function useEquipmentList(condominiumId: string) {
  return useQuery({
    queryKey: ['equipment', condominiumId],
    queryFn: () => equipmentService.findAll(condominiumId),
    staleTime: 1000 * 60 * 5,   // 5 min de cache
  })
}

export function useEquipmentDetail(id: string) {
  return useQuery({
    queryKey: ['equipment', id],
    queryFn: () => equipmentService.findOne(id),
  })
}

export function useMarkInspectionDone() {
  const queryClient = useQueryClient()
  return useMutation({
    mutationFn: inspectionService.markDone,
    onSuccess: () => {
      queryClient.invalidateQueries({ queryKey: ['equipment'] })
      queryClient.invalidateQueries({ queryKey: ['inspections'] })
    },
  })
}
\end{verbatim}

O \texttt{invalidateQueries} após uma mutação garante que o dashboard e a
lista de equipamentos reflitam o novo status imediatamente, sem reload manual.

\subsection{Schemas Zod Compartilhados}

Os schemas de validação em \texttt{lib/schemas.ts} são os mesmos usados
pelo backend para validar os bodies das requisições. Isso garante que o
frontend nunca envie um campo inesperado e que os tipos TypeScript gerados
pelo Zod sejam idênticos nos dois lados.

\begin{verbatim}
// lib/schemas.ts  (compartilhado entre frontend e backend)

export const markInspectionDoneSchema = z.object({
  executedAt:   z.string().datetime(),
  executedById: z.string().uuid(),
  notes:        z.string().optional(),
})

export const addEquipmentSchema = z.object({
  name:            z.string().min(1),
  type:            z.enum(['ELEVATOR','LIGHTNING_ROD','WATER_TANK','OTHER']),
  condominiumId:   z.string().uuid(),
  frequencyDays:   z.number().int().positive(),
  maintenanceType: z.enum(['ROUTINE','PREVENTIVE','CORRECTIVE']),
  description:     z.string().min(1),
})

// Tipos inferidos automaticamente pelo Zod
export type MarkInspectionDoneInput = z.infer<typeof markInspectionDoneSchema>
export type AddEquipmentInput       = z.infer<typeof addEquipmentSchema>
\end{verbatim}

\subsection{Proteção de Rotas por Papel}

O React Router protege as rotas com um componente \texttt{ProtectedRoute}
que verifica o token e o papel antes de renderizar a página:

\begin{verbatim}
// App.tsx

<Routes>
  <Route path="/login"    element={<LoginPage />} />
  <Route path="/register" element={<RegisterPage />} />

  <Route element={<ProtectedRoute requiredRole="MANAGER" />}>
    <Route path="/dashboard"       element={<DashboardPage />} />
    <Route path="/equipment"       element={<EquipmentListPage />} />
    <Route path="/equipment/:id"   element={<EquipmentDetailPage />} />
  </Route>

  <Route element={<ProtectedRoute requiredRole="RESIDENT" />}>
    <Route path="/resident" element={<ResidentPage />} />
  </Route>
</Routes>
\end{verbatim}

\vspace{1cm}
\hrule
\vspace{0.3cm}
{\small\color{gray} Documento vivo — atualizar conforme novas decisões forem tomadas.}

\end{document}